\chapter{Construcción del modelo}\label{cap:construccion}

\begin{entregable}{(a): pasos seguidos para construir el modelo}
Partiendo de \res{n-candidatas} explicativas candidatas y la región censal: (1) se eliminan una a
una las variables con colinealidad severa; (2) se diagnostica la estructura del error del modelo
máximo para elegir cómo hacer inferencia; (3) se seleccionan los efectos principales por
eliminación hacia atrás con contrastes robustos; (4) se reincorporan las confusoras de las
exposiciones de interés; (5) se contrastan interacciones respetando el principio jerárquico; y
(6) se comparan los modelos anidados resultantes. Cada paso es una opción de la configuración y
queda registrado en la corrida.
\end{entregable}

\section{Especificación de partida}

El modelo de partida incluye todas las explicativas que pueden serlo (\cref{cap:datos}), con las
transformaciones logarítmicas de la población, la renta y los ensayos clínicos, y la región como
variable cualitativa con el Sur de referencia. La incidencia se incluye: la mortalidad es, en
buena medida, incidencia por letalidad, y omitirla haría que los coeficientes mezclaran ambas
cosas. Con ella, el resto de coeficientes se interpreta a igual número de diagnósticos por
habitante; el \cref{cap:modelo-final} compara ambos modelos (\cref{tab:incidencia}).

Las explicativas continuas se \emph{centran} en su media. El centrado no cambia el ajuste ni las
pendientes de un modelo sin interacciones, pero hace interpretable la constante (mortalidad
esperada de un condado medio del Sur) y, sobre todo, da sentido a los efectos principales cuando
hay interacciones: cada uno es el efecto de su variable en la región de referencia y con las demás
en su media, y desaparece la colinealidad no esencial entre una variable y su producto.

\section{Colinealidad}\label{sec:colinealidad}

Cuando una explicativa es casi combinación lineal de las demás, su coeficiente se estima con una
varianza enorme y se vuelve inestable. Se mide con el factor de inflación de la varianza
\[
\mathrm{FIV}_j=\frac{1}{1-R_j^2},
\]
donde $R_j^2$ es el coeficiente de determinación de la regresión de $X_j$ sobre las demás
explicativas: la varianza de $\hat\beta_j$ es $\mathrm{FIV}_j$ veces la que tendría con
explicativas incorreladas. Su inversa es la \emph{tolerancia} que imprime SPSS. Complementa al
FIV el número de condición de Belsley \parencite{belsley1980}, que detecta dependencias que
implican a varias variables a la vez.

El FIV máximo del modelo de partida es \res{fiv-max-antes} y el número de condición
\res{ncond-antes}. Se eliminó en cada paso la variable de mayor FIV, recalculando todos los FIV,
hasta que ninguno superó 10 (\cref{tab:poda}): salen \res{podadas}. Las cuatro son redundantes
con otra que permanece: el seguro público incluye al sólo público y el privado es casi su
complementario; la población casada y la proporción de hogares de matrimonios miden casi lo mismo;
y la renta mediana está fuertemente correlacionada, en sentido inverso, con la pobreza
($|r|>0{,}85$ en los cuatro casos). Así, la dimensión de la renta sigue representada en el
modelo máximo por la pobreza.
Tras la poda el FIV máximo es \res{fiv-max-despues} y el número de condición \res{ncond-despues}.

\begin{table}[htbp]
\centering
\small
\caption{Poda iterativa por factor de inflación de la varianza (umbral 10)}\label{tab:poda}
\begin{tabularx}{\linewidth}{r X r X r}
\toprule
Paso & Eliminada & FIV & Más correlada con & $r$ \\
\midrule
1 & Seguro público & \num{25.9} & Solo seguro público & \num{0.867} \\
2 & Población casada & \num{14.6} & Hogares de matrimonios & \num{0.873} \\
3 & Seguro privado & \num{13.9} & Solo seguro público & \num{-0.887} \\
4 & log(Renta mediana) & \num{11.0} & Pobreza & \num{-0.855} \\
\bottomrule
\end{tabularx}
\par\smallskip{\footnotesize\raggedright Número de condición: \num{972.2} antes y \num{187.1} después de la poda.\par}
\end{table}


\figura{fiv}{Factor de inflación de la varianza de cada explicativa antes y después de la poda
(escala logarítmica); la línea vertical marca el umbral de 10.}

\section{Estructura del error y estrategia de inferencia}\label{sec:estructura-error}

Antes de seleccionar variables hay que decidir cómo se calculan los errores típicos, porque de
ellos dependen todos los contrastes de la selección. El modelo máximo
($R^2=\res{r2-inicial}$) presenta dos violaciones claras de las hipótesis clásicas:

\begin{itemize}
  \item \textbf{Heterocedasticidad ligada a la población.} Una tasa calculada con pocas muertes es
  muy ruidosa: si las muertes de un condado de $n$ habitantes se comportan como un proceso de
  Poisson, la varianza de su tasa es proporcional a $1/n$. La varianza de los residuos en el
  decil de condados más pequeños (\res{var-decil1}) es \res{razon-var-deciles} veces la del decil
  más grande (\res{var-decil10}); el contraste de Breusch-Pagan frente a $1/n$ lo confirma
  ($\mathrm{LM}=\res{bp-pob-antes}$, \res{bp-pob-antes-p}). Se estimó la función de varianza
  \[
  \sigma^2(n)=a+\frac{b}{n},\qquad \hat a=\res{vf-a},\quad \hat b=\res{vf-b}\cdot10^{6}\ \text{(\res{vf-b-p})},
  \]
  por mínimos cuadrados ponderados factibles iterados. Su lectura es sustantiva: en un condado de
  \num{10000} habitantes el \res{frac-muestreo-10000}\,\% de la varianza residual es ruido de
  muestreo de la propia tasa; en uno de un millón, el \res{frac-muestreo-1000000}\,\%.
  \item \textbf{Dependencia dentro de cada estado.} Los condados de un mismo estado comparten
  registro de cáncer, sistema sanitario y políticas. La correlación intraclase de los residuos por
  estado es \res{icc-antes} ($F=\res{icc-antes-f}$, \res{icc-antes-p}): los errores no son
  independientes.
\end{itemize}

\figura{varianza}{Varianza residual media por decil de población. Izquierda: modelo inicial y
función de varianza ajustada $a+b/n$. Derecha: modelo final, residuos MCO frente a residuos
ponderados por la función de varianza (MCPF).}

Ninguna de las dos violaciones sesga los coeficientes MCO, pero ambas invalidan sus errores
típicos clásicos, $\hat\sigma^2(X^{\top}X)^{-1}$. Hay dos remedios:
\begin{enumerate}
  \item \textbf{Mantener MCO y corregir la varianza.} El estimador «sándwich»
  \[
  \widehat{\Var}(\hat{\boldsymbol\beta})=(X^{\top}X)^{-1}
  \Bigl(\sum_{g=1}^{G}X_g^{\top}\hat{\mathbf e}_g\hat{\mathbf e}_g^{\top}X_g\Bigr)(X^{\top}X)^{-1}
  \]
  agrupa los residuos por estado ($G$ conglomerados), de modo que es válido con
  heterocedasticidad de cualquier forma y con correlación arbitraria dentro de cada estado
  \parencite{white1980,cameron2015}. Con $G=\res{n-estados-final}$ estados, los contrastes usan la
  $t$ de Student con $G-1=\res{gl-inferencia}$ grados de libertad, como recomiendan
  \textcite{cameron2015} para un número moderado de conglomerados.
  \item \textbf{Ponderar} por $1/\hat\sigma^2(n)$ (MCPF), más eficiente si la función de varianza es
  correcta.
\end{enumerate}

\begin{idea}[Decisión.]
La orientación pide un modelo MCO, de modo que el estimador principal es MCO con errores típicos
robustos por conglomerados de estado. Los errores típicos clásicos se presentan como salida tipo
SPSS (\cref{sec:salida-spss}) y los HC3 \parencite{mackinnon1985} como referencia; la estimación
ponderada por la función de varianza es una especificación de sensibilidad
(\cref{cap:sensibilidad}). Todos los contrastes de la selección usan la covarianza por
conglomerados.
\end{idea}

\section{Selección de variables}

Se aplicó la eliminación hacia atrás: partiendo del modelo máximo (\res{k-maximo} términos,
$n=\res{n-maximo}$ condados con todas las variables), en cada paso se retira el término con mayor
$p$ en el contraste de Wald robusto mientras ese $p$ supere $\alpha=0{,}05$. La región, con tres
indicadoras, se contrasta como un único término (un $F$ con tres grados de libertad). En
\res{n-pasos-seleccion} pasos salen, por este orden, \res{eliminadas} (\cref{tab:seleccion}).

\begin{table}[htbp]
\centering
\small
\caption{Traza de la selección de variables (contrastes de Wald con covarianza robusta por conglomerados)}\label{tab:seleccion}
\begin{tabularx}{\linewidth}{r l X r r}
\toprule
Paso & Acción & Término & $F$ & $p$ \\
\midrule
1 & sale & Pobreza & \num{0.027} & \num{0.869} \\
2 & sale & log(Población) & \num{0.019} & \num{0.890} \\
3 & sale & Población asiática & \num{0.118} & \num{0.733} \\
4 & sale & Sin bachillerato (18–24) & \num{0.112} & \num{0.739} \\
5 & sale & Empleo (16+) & \num{0.993} & \num{0.324} \\
6 & sale & log(1 + ensayos per cápita) & \num{1.125} & \num{0.294} \\
7 & sale & Población negra & \num{1.198} & \num{0.279} \\
8 & sale & Hogares de matrimonios & \num{0.333} & \num{0.566} \\
9 & sale & Grado universitario (18–24) & \num{0.857} & \num{0.359} \\
10 & sale & Bachillerato (25+) & \num{2.465} & \num{0.123} \\
11 & sale & Población nativa o multirracial & \num{3.802} & \num{0.057} \\
12 & sale & Seguro privado vía empleador & \num{1.792} & \num{0.187} \\
\bottomrule
\end{tabularx}

\end{table}


Para comprobar que el resultado no es un artefacto del método se repitió la selección hacia
delante y por pasos: las tres estrategias coinciden en \res{n-comunes-estrategias} términos
(\res{comunes-estrategias}). En cambio, la misma eliminación hacia atrás con errores típicos
\emph{clásicos} retiene \res{n-ingenua} términos, \res{n-ingenua-extra} más
(\res{ingenua-extra}): al ignorar la dependencia por estado y la heterocedasticidad subestima los
errores típicos y declara significativas variables que no lo son.

\section{Confusión}

La selección por significación puede descartar una variable que, sin ser significativa, cambia el
coeficiente de un factor de interés: una \emph{confusora}. Se aplicó el criterio del cambio en la
estimación \parencite{maldonado1993} a las exposiciones socioeconómicas que motivan el estudio
(educación, pobreza y renta): se reincorpora la variable eliminada cuya inclusión cambia más de un
10\,\% el coeficiente de alguna exposición retenida, y se repite hasta que ninguna lo hace. Se
reincorpora \res{confusoras}: su ausencia cambia un \res{confusion-bach25}\,\% el coeficiente de
\res{confusion-bach25-en}. Las dos variables describen la misma distribución educativa (la
proporción con bachillerato y la proporción con título universitario), y estimar el efecto del
título universitario sin controlar por el bachillerato mezclaría ambos escalones educativos.

\section{Interacciones}

Se contrastaron \res{n-interacciones-candidatas} interacciones con fundamento sustantivo: si los
efectos de la incidencia, la educación, la cobertura sólo pública, el desempleo o la pobreza
difieren entre regiones, y si el efecto de la incidencia depende del nivel educativo. Sólo se
contrasta una interacción si sus efectos principales están en el modelo (principio jerárquico), y
en cada ronda los $p$ de todas las candidatas se corrigen por Holm para controlar la tasa de error
de familia. Entran \res{n-interacciones}: \res{interacciones}
(\cref{tab:interacciones}). La primera tiene $F=\res{int-grado25xregion-f}$
(\res{int-grado25xregion-p}; corregido, \res{int-grado25xregion-padj}); la segunda,
$F=\res{int-desempleoxregion-f}$ (\res{int-desempleoxregion-p}; corregido,
\res{int-desempleoxregion-padj}).

\begin{table}[htbp]
\centering
\footnotesize
\caption{Contraste de interacciones candidatas por rondas (corrección de Holm sobre las candidatas de cada ronda)}\label{tab:interacciones}
\begin{tabularx}{\linewidth}{r X r r r r l}
\toprule
Ronda & Interacción & g.\,l. & $F$ & $p$ & $p$ Holm &  \\
\midrule
1 & Incidencia de cáncer $\times$ Región censal & 3 & \num{2.49} & \num{0.072} & \num{0.216} &  \\
1 & Grado universitario (25+) $\times$ Región censal & 3 & \num{7.27} & $<\num{0.001}$ & \num{0.002} & \textbf{entra} \\
1 & Solo seguro público $\times$ Región censal & 3 & \num{1.76} & \num{0.167} & \num{0.216} &  \\
1 & Desempleo (16+) $\times$ Región censal & 3 & \num{3.94} & \num{0.014} & \num{0.055} &  \\
1 & Incidencia de cáncer $\times$ Grado universitario (25+) & 1 & \num{2.89} & \num{0.096} & \num{0.216} &  \\
2 & Incidencia de cáncer $\times$ Región censal & 3 & \num{2.49} & \num{0.071} & \num{0.143} &  \\
2 & Solo seguro público $\times$ Región censal & 3 & \num{1.40} & \num{0.255} & \num{0.255} &  \\
2 & Desempleo (16+) $\times$ Región censal & 3 & \num{4.22} & \num{0.010} & \num{0.040} & \textbf{entra} \\
2 & Incidencia de cáncer $\times$ Grado universitario (25+) & 1 & \num{4.68} & \num{0.036} & \num{0.107} &  \\
3 & Incidencia de cáncer $\times$ Región censal & 3 & \num{2.65} & \num{0.060} & \num{0.120} &  \\
3 & Solo seguro público $\times$ Región censal & 3 & \num{0.22} & \num{0.881} & \num{0.881} &  \\
3 & Incidencia de cáncer $\times$ Grado universitario (25+) & 1 & \num{5.16} & \num{0.028} & \num{0.083} &  \\
\bottomrule
\end{tabularx}

\end{table}


\section{Comparación de modelos anidados}\label{sec:comparacion}

La \cref{tab:comparacion} compara, sobre la misma muestra de \res{n-maximo} condados, el modelo
con sólo la región, el máximo tras la poda, el seleccionado y el final con interacciones. Cada fila
se contrasta con la anterior con el $F$ parcial clásico y con el Wald robusto.

\begin{table}[htbp]
\centering
\footnotesize\setlength{\tabcolsep}{4pt}
\caption{Comparación de modelos anidados (misma muestra y mismos pesos)}\label{tab:comparacion}
\begin{tabularx}{\linewidth}{>{\raggedright\arraybackslash}X r r r r r r r r}
\toprule
Modelo & $k$ & $\bar R^2$ & AIC & BIC & \multicolumn{2}{c}{$F$ clásico} & \multicolumn{2}{c}{Wald robusto} \\ \cmidrule(lr){6-7}\cmidrule(lr){8-9} & & & & & $F$ & $p$ & $F$ & $p$ \\
\midrule
Nulo (región) & 3 & \num{0.154} & \num{25083} & \num{25107} & --- & --- & --- & --- \\
Máximo (tras la poda) & 24 & \num{0.544} & \num{23436} & \num{23584} & \num{110.70} & $<\num{0.001}$ & \num{106.52} & $<\num{0.001}$ \\
Seleccionado (efectos principales) & 13 & \num{0.537} & \num{23462} & \num{23545} & \num{4.35} & $<\num{0.001}$ & \num{2.13} & \num{0.036} \\
Final (con interacciones) & 19 & \num{0.543} & \num{23433} & \num{23551} & \num{6.85} & $<\num{0.001}$ & \num{5.71} & $<\num{0.001}$ \\
\bottomrule
\end{tabularx}
\par\smallskip{\footnotesize\raggedright Cada fila se compara con la anterior (bloque de términos que se añade o se retira). El $F$ clásico supone errores independientes y homocedásticos; el Wald robusto usa la covarianza por conglomerados de estado, la misma que guía la selección.\par}
\end{table}


Tres lecturas:
\begin{itemize}
  \item Las explicativas aportan muchísimo frente a la región sola ($\bar R^2$ pasa de
  \res{cmp-nulo-r2aj} a \res{cmp-maximo-r2aj}).
  \item Los \res{cmp-seleccionado-gl} parámetros que retira la selección no son significativos uno a
  uno, pero conjuntamente el Wald robusto da \res{cmp-seleccionado-p-rob}: aportan, en bloque, una
  mejora pequeña del ajuste ($\bar R^2$ de \res{cmp-seleccionado-r2aj} frente a
  \res{cmp-maximo-r2aj}). El BIC, que penaliza más la complejidad, prefiere el modelo seleccionado
  (\res{cmp-seleccionado-bic} frente a \res{cmp-maximo-bic}); el AIC prefiere el máximo. Se opta por
  el modelo parsimonioso para interpretar, y el modelo máximo se reestima como especificación de
  sensibilidad (\cref{cap:sensibilidad}) para comprobar que las conclusiones no dependen de esta
  elección.
  \item Las interacciones mejoran el ajuste de forma significativa con ambos contrastes
  (Wald robusto \res{cmp-final-p-rob}) y el modelo final tiene el menor AIC de los cuatro.
\end{itemize}

\begin{idea}[Muestra de estimación.]
La selección y la comparación se hacen sobre los \res{n-maximo} condados con todas las candidatas,
para que todos los modelos sean comparables. Una vez fijados los términos, el modelo final se
reestima con todos los condados completos en \emph{sus} variables, \res{n-final}: excluir
condados por variables que el modelo no usa sólo restaría precisión.
\end{idea}
